
\documentclass{configs/idp-model}
% incluir seus pacotes aqui
% obs: cuidado para não conflitar com os pacotes inclusos neste template.



%PREENCHER TODOS OS CAMPOS ABAIXO
\autor{Linus Torvalds}


% Utilizar um dos dois cursos abaixo
\cienciadacomputacao
%\engenhariadesoftware


\titulo{Um título muito interessante para um trabalho}
\ano{2026}

% Usar orientador ou orientadora
% - Argumento opcional caso o orientador seja de uma instituição diferente
% \orientador[IDP]{Dr. Alan Turing}
\orientadora{Dra. Ada Lovelace}

% Para membros da banca, são três argumentos obrigatórios:
% 1 - Nome do membro avaliador
% 2 - Se o mesmo é membro interno ou externo à instituição
% 3 - Instituição a qual o membro está vinculado
\membrobancai{Dr. Dom Quixote de la Mancha}{Internal Member}{IDP}
\membrobancaii{Dra. Anna Karienina}{Internal Member}{IDP}

\palavraschave{LaTeX, metodologia científica, trabalho de conclusão de curso, artigo, tecnologia}
\keywords{{LaTeX, metodologia científica, trabalho de conclusão de curso, paper, technology}}

\dataaprovacao{25/06/2026}

\epigrafeautor{Edsger Wybe Dijkstra}   % opcional



\begin{document}

% Não alterar
\pagenumbering{roman}
\capaprincipal
\folharosto
\folhaaprovacao

% Opcional
\dedicatoria     % opcional
\agradecimentos  % opcional
% se for usar epígrafe, defina também o autor da mensagem com o comando \epigrafeautor
\epigrafe        % opcional


% Não alterar
% Resumo e abstract - alterar os arquivos na pasta "partes" NÃO ALTERAR O NOME DELES
\abstract
\resumo

%List of Figures e Tables são optativos. Caso não esteja utilizando no seu projeto, comentar as linhas abaixo
\listoffigures
\listoftables

\tableofcontents
\pagenumbering{arabic}


% Para as seções, o comando \secao deve ter a forma \secao{título}{arquivo .tex}

% Os comandos \subsecao e \subsubsecao seguem a forma \comando{Título}, sem a necessidade de passar um arquivo .tex

\secao{Introduction}{partes/introduction}
\secao{Literature Review}{partes/literature-review}
\secao{Methodology}{partes/methodology}
\secao{Results and Discussion}{partes/results}
\secao{Conclusion}{partes/conclusion}

% Não alterar
\referencias

% Opcionais
\apendice[ap:meuapendice]{Título do apêndice A}{partes/apendicei}
% Alternativa sem o primeiro parâmetro também funciona
% \apendice{Título do apêndice A}{partes/apendicei}

\anexo{Título do Anexo A}{partes/anexoi}


% Não alterar
\capafinal


\end{document}
